\documentclass{situationist-preprint}

\title{On Holes}
\msid{MET-15-a}
\msfamily{MET-15}
\msdiscipline{Metaphysics}
\mstype{Principal manuscript}
\mskeywords{holes; dependent particulars; immaterial objects; boundaries; absences; ontological dependence}

\begin{document}
\maketitle

\begin{abstract}
There are holes in cheese, in doughnuts and in the ground, and we count, measure and refer to them. Do holes exist? We argue that they do, and that they are dependent immaterial particulars located at the boundaries of their hosts. Holes are particulars because they are countable, located and have shapes and sizes of their own; they are immaterial because they are not made of the matter that surrounds them or of the matter that may fill them; and they are dependent because no hole can exist without a host whose surface bounds it. We reject the view of David and Stephanie Lewis that holes are their material linings and the eliminativist view that hole talk can be paraphrased away. We also argue that holes are not absences. A hole can be filled with water and remain the same hole, whereas an absence of matter cannot be filled without ceasing to be; and the identity of a hole is fixed by its host, whereas an absence has no host. The result is therefore consistent with the denial, in MET-16, that there are absences.
\end{abstract}

\section{The problem}

Holes are puzzling because they seem to be made of nothing. A hole in a piece of cheese has a shape, a size and a location; it can be enlarged, filled and counted; but there is nothing at the place where it is, or at least nothing that belongs to the cheese. Ordinary language quantifies over holes freely (\enquote{there are three holes in this sock}), and the quantification supports inference and counting. If we take it at face value, we must say what holes are.

\citet{LewisLewis1970} set out the options in a dialogue between Argle, who denies that there are immaterial objects, and Bargle, who presses him with holes. Argle first tries to paraphrase hole talk into talk of perforated objects, and then concludes that holes exist but are material: a hole is the hole-lining, the part of the host that surrounds the cavity. \citet{CasatiVarzi1994} developed the alternative on which holes are immaterial bodies that depend on their hosts, and placed them in a general theory of parasitic or \enquote{superficial} entities. A third option, defended in the debate by eliminativists, is that there are no holes, only holed things: the sentence \enquote{there is a hole in the cheese} is true because the cheese is perforated, and it does not commit us to holes. We argue for a version of Casati and Varzi's view, and against both the lining view and eliminativism.

\section{Framework}

\begin{definition}\label{d:host}
A \emph{host} is a material object whose surface has a configuration (a depression, a cavity or a perforation) that bounds a region not occupied by the host.
\end{definition}

\begin{definition}\label{d:hole}
A \emph{hole} in a host $h$ is a particular whose location is the region bounded by such a configuration of $h$'s surface, and whose existence and shape depend on that configuration.
\end{definition}

\begin{principle}[Ontological commitment]\label{p:commit}
If a body of discourse quantifies over, counts and refers to items of a kind, supports inferences that depend on their identity and number, and admits no paraphrase that preserves these inferences without quantifying over them, then items of that kind exist.
\end{principle}

\begin{principle}[Indiscernibility of matter]\label{p:matter}
If $x$ is made of some matter, then $x$ has the properties that being made of that matter confers: it is where that matter is, and it changes when that matter is removed or replaced.
\end{principle}

\section{The argument}

\begin{lemma}\label{l:particular}
Hole discourse satisfies the antecedent of Principle~\ref{p:commit}.
\end{lemma}

\begin{proof}
We count holes (\enquote{there are as many holes in this slice as in that one}), refer to them (\enquote{the hole you drilled yesterday has widened}) and draw inferences that depend on their identity (\enquote{the hole that was filled with water this morning is the hole the mouse made}). The paraphrase into talk of perforated hosts preserves some of these inferences but not all. A host that is \enquote{doubly perforated} has two holes; but the paraphrase must then introduce a predicate for each number of holes, and it cannot express that the hole in one host is larger than the hole in another, or that two hosts share a hole, as when a hole runs through two adjoining boards, without reintroducing a variable that ranges over holes \citep{LewisLewis1970,CasatiVarzi1994}.
\end{proof}

\begin{theorem}\label{t:exist}
Holes exist, and they are particulars.
\end{theorem}

\begin{proof}
By Lemma~\ref{l:particular} and Principle~\ref{p:commit}, holes exist. They are particulars because they are counted, located and individually referred to, and because each has a shape and a size that may differ from those of other holes.
\end{proof}

\begin{proposition}\label{p:immaterial}
Holes are immaterial.
\end{proposition}

\begin{proof}
A hole is not made of its host's matter, since it is located where the host's matter is not. The lining view identifies the hole with a part of the host that surrounds the cavity. But the lining is where the host's matter is, and the hole is not; the lining is made of cheese and can be eaten, and the hole cannot. By Principle~\ref{p:matter}, the hole is not the lining. Nor is a hole made of whatever fills it. A hole in the ground may fill with rainwater and drain again; it is the same hole throughout, while the water comes and goes. By Principle~\ref{p:matter}, a thing made of the water would change when the water was replaced, and the hole does not. A hole is therefore made of no matter at all.
\end{proof}

\begin{proposition}\label{p:dependent}
Every hole depends for its existence on a host, and is located at the boundary of that host.
\end{proposition}

\begin{proof}
A hole in a piece of cheese ceases to exist if the cheese is melted down, although nothing happens at the region the hole occupied except that the surrounding cheese flows into it or away from it. Its existence depends on the configuration of the host's surface, which is what Definition~\ref{d:hole} records. Its location is the region that the configuration bounds, and its shape is the shape of that boundary. There is no hole that is not bounded by some host's surface, so every hole has a host.
\end{proof}

\begin{corollary}\label{c:main}
Holes are dependent immaterial particulars located at the boundaries of their hosts.
\end{corollary}

\begin{proof}
By Theorem~\ref{t:exist} and Propositions~\ref{p:immaterial} and~\ref{p:dependent}.
\end{proof}

\section{Holes are not absences}

It is tempting to say that a hole is an absence of matter, and so to treat the theory of holes as a branch of the theory of absences. We reject this, for three reasons, and the rejection matters because MET-16 argues that there are no absences, while we have argued that there are holes.

First, holes can be filled without ceasing to exist. A hole in the ground full of water is still a hole, and a mechanic who fills a hole in a panel with putty may later drill the putty out and find the hole as it was. An absence of matter, by contrast, cannot be filled: the absence of water from a region is precisely what ceases when water arrives there. If holes were absences of matter, a filled hole would be a contradiction, and it is not.

Secondly, holes have hosts and absences do not. The identity of a hole is fixed by the host whose surface bounds it: the hole in this slice of cheese is not the hole in that one, even if the two cavities were congruent and the slices were swapped. An absence of matter from a region has no host; it is individuated, if at all, by the region and by what is missing from it. The empty space between two stars is a region from which matter is absent, but it is not a hole in anything.

Thirdly, the truths that we might express by speaking of absences are made true by what exists and fails to exist, as MET-16 argues, and they require no entity. The truths about holes are not of this kind. That a sock has three holes is not made true by the non-existence of anything; it is made true by the configuration of the sock's surface, and by the three particulars that this configuration bounds. Truths about holes are positive truths about dependent particulars. Truths about absences are negative truths, which require no particular at all.

The two theses are therefore consistent. There are holes, which are dependent immaterial particulars; there are no absences; and the region a hole occupies is, if the hole is unfilled, also a region from which matter is absent, but this coincidence of location does not make the hole an absence, any more than a shadow is an absence of light because there is no light where the shadow is.

\section{Objections}

\begin{objection}[Parsimony]
Immaterial particulars are an ontological extravagance; the lining view gets by with material objects alone.
\end{objection}

\begin{reply}
The lining view gets by with material objects at the cost of misdescribing holes. It makes a hole edible, gives it the mass of the surrounding cheese, and makes the size of the hole depend on how thick a lining one chooses. Parsimony favors the more economical of two theories that describe the phenomena equally well, and the lining view does not describe them equally well. Moreover, holes add no fundamental kind: they are fixed by the surfaces of their hosts, and their existence is grounded in the configurations of material objects.
\end{reply}

\begin{objection}[Location]
An immaterial particular with a location is a strange thing. What occupies the region where an unfilled hole is?
\end{objection}

\begin{reply}
The hole does. That is what it is for the hole to be located there. The objection assumes that occupying a region requires being material, and the case of holes shows that it does not. Holes are not alone in this: boundaries, surfaces and shadows are located and immaterial, and on the account of \citet{CasatiVarzi1994} holes are one species of a broader genus of dependent superficial entities.
\end{reply}

\begin{objection}[Composition]
If the hole is immaterial, it has no parts; but we speak of half a hole and of the deeper part of a hole.
\end{objection}

\begin{reply}
A hole has parts in the sense in which a region does: the region it occupies has subregions, and the deeper part of the hole is the part located at the deeper subregion. The question when some things compose something, which MET-06 answers for material objects, does not arise for holes in the same form, since the parts of a hole are not independent things that come together, but divisions of a particular that exists in virtue of its host's surface.
\end{reply}

\section{Conclusion}

Holes exist. They are particulars, since we count, measure and refer to them in ways that resist paraphrase; they are immaterial, since they are neither their linings nor whatever fills them; and they are dependent, since they exist in virtue of the surface configurations of their hosts and are located at the regions those configurations bound. They are not absences: they can be filled, they have hosts, and truths about them are positive truths. The existence of holes is therefore no counterexample to the view that there are no absences.

\begin{thebibliography}{9}\small
\bibitem[Casati and Varzi(1994)]{CasatiVarzi1994} Casati, R. and Varzi, A.~C. (1994). \emph{Holes and Other Superficialities}. MIT Press.
\bibitem[Lewis and Lewis(1970)]{LewisLewis1970} Lewis, D. and Lewis, S. (1970). Holes. \emph{Australasian Journal of Philosophy}, 48(2), 206--212.
\end{thebibliography}

\end{document}
